\section{Methods}
\label{sec:methods}

\subsection{Model overview}

We constructed a state-transition (Markov) model to compare lifetime outcomes
for two hypothetical cohorts of otherwise identical individuals: one cohort that
donates a kidney and one that does not.
Both cohorts enter the model at the age of donation in full health, and the
simulation advances in annual cycles until age~100.
The analysis was conducted from the perspective of the individual donor rather
than the health-care system; outcomes are reported as undiscounted remaining
life-years, consistent with the standard approach for life-expectancy analyses
\cite{gold_cost_1996}.
The base-case age at donation is~40 years; subgroup analyses span ages 25--55.
Each arm was simulated with $5 \times 10^6$ individuals (Monte Carlo) and
cross-validated against a deterministic analytic cohort implementation.

\subsection{Health states}

The model contains seven mutually exclusive health states
(Figure~\ref{fig:markov}):
\textit{Healthy}, \textit{ESRD onset}, \textit{Dialysis},
\textit{Standard waitlist}, \textit{Priority waitlist},
\textit{Post-transplant}, and \textit{Dead} (absorbing).
At ESRD onset the cohort branches: a fraction is listed preemptively
(bypassing dialysis entirely) and the remainder initiates haemodialysis.
The donor arm routes ESRD survivors to the \textit{Priority waitlist},
reflecting the allocation priority afforded to prior living donors (PLD)
under KAS250; the non-donor arm routes to the \textit{Standard waitlist}.
Graft failure returns patients to the \textit{Dialysis} state, bypassing
the ESRD-onset branch.

\subsection{ESRD risk after donation}

Fifteen-year cumulative ESRD incidence was taken from Muzaale et al.\
\cite{muzaale_risk_2014}:
30.8 per 10,000 (95\%~CI 24.3--38.5) for donors and 3.9 per 10,000
(95\%~CI 0.8--8.9) for age-, sex-, and health-matched controls
(8-fold relative risk).
Race-specific donor rates were 74.7/10,000 (Black) and 22.7/10,000
(White) \cite{muzaale_risk_2014}.
The White non-donor baseline was taken from Grams et al.\
\cite{grams_kidney-failure_2016} (sex-averaged 0.05\%/15~yr), because
zero ESRD events were recorded in White matched controls in Muzaale et al.

ESRD hazard was modelled with a Weibull distribution (shape $k = 1.5$),
calibrated via a competing-risk bisection algorithm to reproduce the
15-year cumulative incidence while accounting for background-mortality
attrition in the at-risk pool.
The accelerating hazard ($k > 1$) is supported by the Massie et al.\ 20-year
cumulative incidence of 34 per 10,000 (vs.\ $\approx$21/10,000 predicted
under a constant-hazard model) \cite{massie_quantifying_2017}.
Within-donor subgroup hazard ratios were drawn from Massie et al.:
2.96 (95\%~CI 2.25--3.89) for Black race, 1.88 (95\%~CI 1.50--2.35) for
male sex, and 1.40 per decade (non-Black donors) \cite{massie_quantifying_2017}.

\subsection{Dialysis survival and waitlist listing}

First-year haemodialysis mortality was 22\% and subsequent-year mortality 17\%
($\approx$42\% five-year survival), consistent with USRDS estimates
\cite{usrds_2024}.

\subsubsection*{Preemptive listing}

At ESRD onset, a fraction of patients is listed preemptively before dialysis
initiation, avoiding dialysis mortality entirely.
Preemptive listing probabilities were drawn from USRDS 2025 ADR Figure~7.13
\cite{usrds_2025}: 5.8\% overall for the non-donor arm and 9.4\% for the donor
arm (age 18--44 cohort, reflecting the younger and healthier donor-like
population).

\subsubsection*{Post-dialysis listing probability}

Among patients on dialysis, the conditional annual probability of joining the
waitlist was set to 15\%/year (base case), derived from USRDS 2025 ADR
Figure~7.17 \cite{usrds_2025}.
Approximately 25\% of 18--44 year-old ESRD patients are listed or transplanted
within one year, of whom $\approx$9.4\% were preemptively listed, yielding an
estimated post-dialysis listing rate of $\approx$15.6\%/year for the
donor-like cohort (conservative base case: 15\%/year; one-way sensitivity range
5--30\%/year).

\subsection{Transplant waiting time and waitlist outcomes}

Transplant probability while listed was modelled as a time-homogeneous
exponential process parameterised by the median waiting time, a standard
approximation for waiting-time analyses \cite{krahn_primer_1997}.
For the standard waitlist, the post-KAS250 national median of 985~days
($\approx$32.8~months) was used \cite{schold_effect_2023};
for the PLD priority waitlist, the post-KAS median of 102.6~days was
used \cite{wainright_impact_2017}.
Waitlist mortality was 5.0 deaths per 100 patient-years (SRTR 2023 ADR
Figure~KI~24) and competing removal 6.81\%/year (SRTR 2023 ADR
Figure~KI~22) \cite{srtr_2023}.

\subsection{Post-transplant survival}

Post-transplant annual mortality was derived from age-stratified 5-year
Kaplan--Meier estimates for DDKT recipients in the SRTR 2023 ADR
(Figure~KI~70, 2016--2018 cohort) \cite{srtr_2023}:
0.9\%/year (ages 18--34), 1.8\%/year (35--49), 3.9\%/year (50--64),
and 6.9\%/year ($\geq 65$).
The base case uses DDKT survival because prior living donors who develop
ESRD require a deceased-donor organ.
A one-way sensitivity analysis substitutes LDKT patient survival
(SRTR 2023 ADR Figure~KI~76: 0.4\%, 0.8\%, 1.7\%, and 3.9\%/year by
age band) to bound the effect of this assumption.
Annual graft failure was 2.5\%/year post-year~1 \cite{srtr_2023};
graft failure returned patients to the Dialysis state.

\subsection{Background mortality and donor all-cause survival}

Background mortality in the Healthy state was taken from the 2021 US
Life Tables (CDC, National Vital Statistics Reports Vol.~72 No.~12)
\cite{cdc_2023}, with sex-specific tables used in sex-stratified analyses.
The base-case donor all-cause mortality hazard ratio versus matched
non-donors is 1.0, consistent with US-based studies
\cite{muzaale_risk_2014, segev_perioperative_2010} and the Grams et al.\
2018 meta-analysis (pooled HR 0.984, 95\%~CI 0.743--1.302;
9 studies, $n = 118{,}426$ donors) \cite{grams_mid-_2018}.
An HR of 1.30 (Mjøen et al.\ \cite{mjoeen_long-term_2014}, a Norwegian
cohort in which 80\% of donors were first-degree relatives of the
recipient) was examined as the pessimistic sensitivity scenario.

\subsection{Subgroup analyses}

\subsubsection*{Age at donation}

Analyses were repeated at ages 25, 35, 40, 45, and 55~years at donation.
Donor ESRD rates were scaled by the Massie~2017 age HR (1.40/decade,
non-Black donors) \cite{massie_quantifying_2017}; non-donor baselines
were held at the Muzaale overall matched-control rate.

\subsubsection*{Race}

Race-specific donor ESRD rates (Black 74.7/10,000; White 22.7/10,000)
and race-specific non-donor baselines (Black sex-averaged 0.195\%/15~yr;
White sex-averaged 0.05\%/15~yr) were taken from Muzaale et al.\ and
Grams et al., respectively \cite{muzaale_risk_2014, grams_kidney-failure_2016}.
Race-specific waitlist mortality and post-transplant survival were from
SRTR 2023 Figures~KI~25 and KI~71 \cite{srtr_2023}.

\subsubsection*{Sex}

Sex-specific donor ESRD rates were derived by applying the Massie~2017
male-sex HR (1.88) to the overall donor rate, weighted by the observed
donor sex mix (60\% female; SRTR 2023 ADR Figure~KI~7) \cite{massie_quantifying_2017, srtr_2023}.
For non-donor baselines, we used the Grams~2016 White sex-specific
rates (female 0.04\%/15~yr; male 0.06\%/15~yr) rather than applying
the within-donor sex HR to the control arm, because Massie~2017
hazard ratios describe variation within the donor cohort and are not
appropriate for scaling a different reference population
\cite{grams_kidney-failure_2016}.
This sourcing rule---Muzaale matched-control rates for overall analyses;
Grams direct rates wherever sex or race is stratified---was applied
consistently across all matrix subgroup figures (see Supplementary Methods).

\subsection{Sensitivity and uncertainty analyses}

\subsubsection*{Probabilistic sensitivity analysis}

We conducted a PSA with 200 parameter draws.
Sampled quantities and their distributions were:
ESRD 15-year cumulative risks (beta distributions parameterised from the
Muzaale~2014 published 95\%~CIs);
donor all-cause mortality HR (log-normal, $\mu = {-0.016}$,
$\sigma = 0.143$, derived from the Grams~2018 meta-analysis CI);
Weibull shape $k$ (log-normal, $\sigma = 0.13$, representing structural
uncertainty absent published CIs for this parameter); and
waitlist listing probability (beta, fractional SE~=~40\%, consistent with
the 0.05--0.30 plausible range).
Registry-derived parameters (SRTR, USRDS) have negligible sampling error
relative to structural uncertainty and were fixed at base-case values;
scenario variation in these parameters is captured by the one-way
sensitivity analysis.
We report the mean $\Delta$LE, 95\%~credible interval, and the
probability that donation is life-expectancy--neutral or beneficial.

\subsubsection*{One-way sensitivity analysis}

Each parameter was varied individually while others were held at
base-case values (Table~\ref{tab:owsa}).
Key scenarios included: no priority access (PLD wait set to the
standard 985-day median), donor all-cause mortality HR~=~1.30
(Mjøen pessimistic estimate), ESRD relative risk
$\times 11$ (Mjøen upper-bound ESRD HR), PLD wait 50~days (optimistic),
post-transplant survival at LDKT quality (SRTR Figure~KI~76), and
dialysis mortality~$\pm 50\%$.

\subsection{Calibration}

The model's transplant probability was validated against the SRTR 2023
3-year waitlist outcome distribution (Figure~KI~22, 2019--2021 listing
cohort) \cite{srtr_2023}.
This comparison is non-circular because the transplant probability was
derived from the median waiting time, not from the 3-year outcome
distribution.
The model predicted 44.3\% transplanted at three years versus 44.3\%
observed (DDKT + LDKT combined; $\Delta = 0\%$).
The waitlist removal rate was back-calculated via a competing-risk
bisection algorithm to reproduce the SRTR-observed 19.1\% removed at
three years (12.6\%/year); the na\"{i}ve formula $1-(1-\mathrm{CIF})^{1/3}$,
which ignores competing depletion by transplantation and death, was not
used.
The model over-predicted three-year waitlist mortality (10.0\% vs.\
6.8\% observed, $\Delta = +47\%$), likely reflecting the application of a
2023 cross-sectional mortality rate (confirmed 5.0/100 patient-years from
the raw SRTR data file) to a 2019--2021 listing cohort; this bias is
symmetric across both arms and conservative (slightly overstates the cost
of the waitlist period for both donors and non-donors equally).
No independent cohort exists tracing the complete pathway from living
donation through ESRD onset to transplantation; these checks represent
the available face-validity evidence for the model.

\subsection{Analytic cycle length and discounting}

Annual (one-year) cycles were used throughout.
Life-years accumulated with an end-of-cycle convention; the resulting
systematic underestimate of approximately 0.5~years applies equally to
both arms and cancels in $\Delta$LE.
Future life-years were not discounted, consistent with the convention
for pure life-expectancy analyses \cite{gold_cost_1996}.

\subsection{Software and reproducibility}

All analyses were implemented in Python~3.10 (NumPy, pandas, matplotlib).
The simulation, parameter assembly pipeline, sensitivity analyses, and
calibration checks are fully reproducible from the public code repository
at \url{https://github.com/rupumped/kidney-donor-le}.
